\appendix

\section{Hardware and software}
\label{sec:appendix-hardware}

\begin{widetable}
\caption{The two devices. Core counts and memory bandwidth are the
manufacturer's specifications; memory capacity, compute capability and
software versions were read from the devices.}
\label{tab:hardware}
\footnotesize
\begin{tabular}{p{0.2\linewidth}p{0.36\linewidth}p{0.36\linewidth}}
\toprule
 & Jetson Orin NX 16\,GB & Jetson AGX Thor \\
\midrule
GPU & NVIDIA Ampere, 1\,024 CUDA cores, 32 tensor cores & NVIDIA Blackwell, 2\,560 CUDA cores, 96 tensor cores \\
compute capability & 8.7 & 11.0 \\
CPU & 8-core Arm Cortex-A78AE & 14-core Arm Neoverse-V3AE \\
memory & 16\,GB LPDDR5, shared by CPU and GPU & 128\,GB LPDDR5X, shared by CPU and GPU \\
memory bandwidth & 102\,GB/s & 273\,GB/s \\
system software & L4T R36.4.3, CUDA 12.6 & L4T R39.2.1, CUDA 13.2 \\
Rust toolchain & nightly-2026-04-03 & nightly-2026-10-11 \\
cuda-oxide & revision b7fe543 & revision e027ac3 \\
\bottomrule
\end{tabular}
\end{widetable}

Both devices were shared with other work during our experiments. In
particular, language-model services occupied most of the Thor's memory until
the larger training run began, at which point they were stopped.

\FloatBarrier
\section{Benchmarks}
\label{sec:appendix-bench}

All figures in this appendix were measured on the Orin NX with
\code{staleynet bench}. Each timing is the mean of five runs after one
warm-up run, and each table entry is the median of three such measurements,
because the clock frequency of the device varies noticeably from run to run.

\begin{widetable}
\caption{Matrix products at the sizes used by the first model, in GFLOP/s.}
\label{tab:bench-gemm}
\small
\begin{tabular}{lrrrr}
\toprule
product ($M \times N \times K$) & direct & tiled & register-blocked & blocked / tiled \\
\midrule
attention projections ($2048 \times 768 \times 256$) & 46.8 & 258.6 & 579.0 & 2.2 \\
feed-forward ($2048 \times 1024 \times 256$) & 55.3 & 255.5 & 450.6 & 1.8 \\
output heads ($2048 \times 4096 \times 256$) & 31.2 & 222.1 & 458.5 & 2.1 \\
spectrogram ($4096 \times 513 \times 1024$) & 47.7 & 236.1 & 541.7 & 2.3 \\
\bottomrule
\end{tabular}
\end{widetable}

\begin{table}[htbp]
\centering
\caption{One training step of the first model on $4 \times 512$ positions,
forward and backward.}
\label{tab:bench-train}
\small
\begin{tabular}{lr}
\toprule
matrix-product kernel & time per step (ms) \\
\midrule
direct & 2\,930 \\
tiled & 576 \\
register-blocked & 315 \\
\bottomrule
\end{tabular}
\end{table}

\begin{table}[htbp]
\centering
\caption{Other measurements on the Orin NX.}
\label{tab:bench-other}
\small
\begin{tabular}{lr}
\toprule
measurement & result \\
\midrule
log-mel spectrogram of 60\,s, GPU & 47\,ms \\
log-mel spectrogram of 60\,s, CPU (rustfft, 8 threads) & 73\,ms \\
decoding step, key/value cache & 6.2\,ms \\
decoding step, recomputing the context & 17.1\,ms \\
\bottomrule
\end{tabular}
\end{table}

On the Thor, measured during training runs of identical shape, register
blocking reduced a training step of the first model on $4 \times 768$ positions
from 238\,ms to 147\,ms, and a step of the larger model on $8 \times 1\,408$
positions takes about 2.0 seconds.

\FloatBarrier
\section{Tokenizer study}
\label{sec:appendix-tokenizer}

\Cref{tab:tokenizer-study} lists the mean reconstruction error over all 39
held-out songs for the tokenizers we compared, each trained on the same
training songs with 1\,024 prototypes per codebook. The columns give the
number of frames per step, the number of band groups, the number of residual
levels in each group, and the resulting number of codes per step.

\begin{table}[htbp]
\centering
\caption{Held-out reconstruction error of different tokenizers, in log units
per mel band (one log unit is a factor of $e$ in energy, about 4.3\,dB). A frame
covers 10.7\,ms of audio; every codebook has 1\,024 prototypes.}
\label{tab:tokenizer-study}
\small
\begin{tabular}{rrrrr}
\toprule
frames & groups & levels & codes & error (log units per band) \\
\midrule
2 & 1 & 4 & 4 & 0.776 \\
2 & 1 & 8 & 8 & 0.643 \\
1 & 1 & 4 & 4 & 0.745 \\
1 & 1 & 8 & 8 & 0.588 \\
1 & 4 & 2 & 8 & 0.597 \\
1 & 8 & 1 & 8 & 0.612 \\
1 & 1 & 16 & 16 & 0.427 \\
1 & 8 & 2 & 16 & 0.403 \\
1 & 4 & 4 & 16 & 0.401 \\
\bottomrule
\end{tabular}
\end{table}

At 8 codes per step, dividing the bands into groups did not help. With 16
codes per step it helped a little: 16 levels without groups reached 0.427,
and 4 groups of 4 levels reached 0.401. Most of the improvement over the first
tokenizer therefore comes from describing each step with more codes, and a
smaller part, about 6\,\%, from quantizing the band groups separately.

\FloatBarrier
